
\documentclass[12pt,a4paper,oneside]{extarticle}
\usepackage{fontspec}
\setmainfont{Times New Roman}

\usepackage{tikz}
\usetikzlibrary{positioning,arrows.meta}

\usepackage{float}
\usepackage{adjustbox}

\fontsize{13pt}{15pt}\selectfont

\usepackage[vietnamese]{babel}
% Font chính: ưu tiên Times New Roman (theo chuẩn báo cáo).
% Nếu không có, tự động fallback sang TeX Gyre Termes (kèm sẵn TeX Live).
\IfFontExistsTF{Times New Roman}
  {\setmainfont{Times New Roman}}
  {\setmainfont{TeX Gyre Termes}}

%---- Lề trang & line spacing -----------------------------------------------
\usepackage[a4paper, top=2.5cm, bottom=2.5cm, left=3cm, right=2cm]{geometry}
\usepackage{setspace}
\onehalfspacing                            % cách dòng 1.5 (chuẩn báo cáo)

%---- Toán học --------------------------------------------------------------
\usepackage{amsmath, amssymb, amsfonts, amsthm}
\usepackage{bm}                            % cho \bm{...} bold math
\usepackage{mathtools}

%---- Đồ họa, bảng, code ----------------------------------------------------
\usepackage{graphicx}
\usepackage{float}                         % buộc figure đứng đúng chỗ với [H]
\usepackage{caption}
\usepackage{subcaption}
\usepackage{array, tabularx, booktabs, multirow, makecell}
\usepackage{xcolor}
\usepackage{listings}
\usepackage{algorithm}
\usepackage{algpseudocode}
% Custom column type
\newcolumntype{Y}{>{\raggedright\arraybackslash}X}
\newcolumntype{C}[1]{>{\centering\arraybackslash}m{#1}}


%---- Hyperlink & mục lục ---------------------------------------------------
\usepackage[colorlinks=true, linkcolor=black, citecolor=blue,
            urlcolor=blue, bookmarks=true]{hyperref}
\usepackage{tocloft}                       % tùy biến mục lục
\renewcommand{\cfttoctitlefont}{\Large\bfseries}

%---- Header / footer -------------------------------------------------------
\usepackage{fancyhdr}
\pagestyle{fancy}
\fancyhf{}
\fancyfoot[C]{\thepage}
\renewcommand{\headrulewidth}{0pt}

%---- Tiêu đề mục (override) ------------------------------------------------
\usepackage{titlesec}
\titleformat{\section}{\Large\bfseries}{\thesection.}{0.5em}{}
\titleformat{\subsection}{\large\bfseries}{\thesubsection.}{0.5em}{}
\titleformat{\subsubsection}{\normalsize\bfseries}{\thesubsubsection.}{0.5em}{}
\titlespacing*{\section}{0pt}{14pt}{8pt}
\titlespacing*{\subsection}{0pt}{12pt}{6pt}

%---- Cấu hình listings cho code Python -------------------------------------
\definecolor{codegreen}{rgb}{0,0.6,0}
\definecolor{codegray}{rgb}{0.5,0.5,0.5}
\definecolor{codepurple}{rgb}{0.58,0,0.82}
\definecolor{backcolour}{rgb}{0.96,0.96,0.96}
\lstdefinestyle{pystyle}{
    backgroundcolor=\color{backcolour},
    commentstyle=\color{codegreen},
    keywordstyle=\color{magenta},
    numberstyle=\tiny\color{codegray},
    stringstyle=\color{codepurple},
    basicstyle=\ttfamily\footnotesize,
    breakatwhitespace=false,
    breaklines=true,
    captionpos=b,
    keepspaces=true,
    numbers=left,
    numbersep=5pt,
    showspaces=false,
    showstringspaces=false,
    showtabs=false,
    tabsize=2,
    language=Python
}
\lstset{style=pystyle}

%---- Lệnh tiện ích ---------------------------------------------------------
% Theo yêu cầu thầy: chữ thẳng = fact từ paper; chữ nghiêng = nhận xét cá nhân.
% Dùng lệnh \fact{...} và \remark{...} để phân biệt rõ.
\newcommand{\fact}[1]{#1}                  % chữ thẳng - fact trích từ paper
\newcommand{\remark}[1]{\textit{#1}}       % chữ nghiêng - nhận xét cá nhân

% Highlight cho trích dẫn/công thức quan trọng
\newcommand{\important}[1]{\textbf{#1}}

%==============================================================================
\begin{document}

%---- TRANG BÌA -------------------------------------------------------------
\begin{titlepage}
\begin{center}

{\large\bfseries ĐẠI HỌC QUỐC GIA TP HỒ CHÍ MINH}\\[2pt]
{\large\bfseries TRƯỜNG ĐẠI HỌC KHOA HỌC TỰ NHIÊN}\\[2pt]
{\large\bfseries KHOA CÔNG NGHỆ THÔNG TIN}\\[0.5cm]

\includegraphics[width=0.3\textwidth]{figures/hcmus-logo.png}\\[0.5cm]
% \vspace{0cm}

% Logo trường (nếu có)
% \includegraphics[width=0.25\textwidth]{logo_hcmus.png}\\[1cm]

{\large BÁO CÁO MÔN HỌC SÂU}\\[6pt]

% \vspace{1cm}

{\Large\bfseries Đề tài:}\\[8pt]
{\LARGE\bfseries Mô Hình Ngôn Ngữ Lớn Đa Phương Thức Song Ngữ Có Thể Diễn Giải cho Các Tác Vụ Y Sinh Đa Dạng}\\[1.2cm]
{\large\bfseries Dựa trên bài báo:}\\[2pt]
{\large\itshape ``Interpretable Bilingual Multimodal Large Language Model\\
for Diverse Biomedical Tasks''}\\[2pt]
{\large Lehan Wang, Haonan Wang, Honglong Yang, Jiaji Mao,\\
Zehong Yang, Jun Shen, and Xiaomeng Li}\\[2pt]
{\normalsize The Thirteenth International Conference on Learning Representations\\
(ICLR 2025)}\\[1cm]

\begin{flushleft}
\hspace{2cm}\begin{tabular}{ll}
\textbf{GIÁO VIÊN HƯỚNG DẪN:}   & TS. Nguyễn Tiến Huy \\
                                & TS. Lê Thanh Tùng\\[8pt]
                                
\textbf{HỌC VIÊN THỰC HIỆN:}  & 25C15029 -- Lê Thị Tường Vy\\
                              & 25C15008 -- Phạm Ngọc Hiếu\\
                              & 25C15011 -- Nguyễn Huy Hoàng\\
                              & 25C11050 -- Nguyễn Đình Lộc\\
                              & 25C15003 -- Ngô Trương Minh Đạt\\
\end{tabular}
\end{flushleft}

\vfill
{\bfseries TPHCM -- 06/2026}

\end{center}
\end{titlepage}

%---- MỤC LỤC ---------------------------------------------------------------
\renewcommand{\contentsname}{MỤC LỤC}
\pagenumbering{roman}
\tableofcontents
\newpage

% (Tùy chọn) Danh sách hình & bảng
% \listoffigures
% \listoftables
% \newpage

\pagenumbering{arabic}
\setcounter{page}{1}


\newpage


% ============ TÓM TẮT ============
\input{contents/Trong-summary}
 
\newpage
 
% ============ 0. GIỚI THIỆU ============



% ============ 1. NHỮNG NGHIÊN CỨU GẦN ĐÂY ============



% ============ 4. KẾT QUẢ THỰC NGHIỆM ============
\section{Thực nghiệm}


%=============== KẾT LUẬN ===============



%==============================================================================
%  TÀI LIỆU THAM KHẢO
%==============================================================================
% \renewcommand{\refname}{Tài liệu tham khảo}
%% Numbered bibliography style
\bibliographystyle{elsarticle-num}

\input{ref.tex}

\newpage

%==============================================================================
%  PHỤ LỤC
%==============================================================================
\appendix

\section{Phụ lục: Code chương trình}

Code chương trình không có, tác giả không công bố công khai.\\
Link Source code của tác giả: 

\url{https://colab.research.google.com/drive/1Gz0L8AnT4ijrUchrhEXXsnaacrFdenn}

\end{document}
